\section{Related work}

\paragraph{Functional-connectome prediction.}
ABIDE-I aggregated resting-state fMRI and phenotypic data across international
sites to support large-scale autism research \citep{dimartino2014abide}. The
Preprocessed Connectomes Project subsequently released derivatives from
several pipelines, including C-PAC \citep{abidepcp}. Multi-site prediction is
challenging because acquisition, motion, demographics, and preprocessing can
vary across sites. Held-out-site evaluation therefore addresses a stricter
generalization problem than random participant splitting.

\paragraph{Graph aggregation.}
GCNs apply normalized neighborhood aggregation followed by learned feature
transformation \citep{kipf2017gcn}. Repeated diffusion can contract node
representations, although the practical effect depends on architecture,
features, graph construction, and task. In functional connectomes, the problem
is especially relevant because connectivity-row node features already contain
whole-graph information. The benchmark of \citet{han2026rethinking} motivates
testing propagation against a no-propagation control rather than comparing
only different GNN families.

\paragraph{Learned transport.}
Neural sheaf diffusion assigns vector spaces and learned maps to graph
incidences, allowing non-trivial transport before aggregation
\citep{bodnar2022sheaf}. BuNN specializes this idea to flat vector bundles with
orthogonal node maps and diffusion-based propagation
\citep{bamberger2025bundle}. Raw embeddings from different node frames cannot
be compared naively; meaningful collapse diagnostics require transport to a
common frame or a gauge-invariant measure. The present experiment therefore
separates raw map-generator capacity from inter-node transport using a
learned-local control.
